% ══════════════════════════════════════════════════════════════════════════════
%  ADVANCED FORCES — the above-level book (from 2027)
%  Faculty decision, Nov 2026: every student completes the at-level lesson books;
%  students who excel move on to this book. Its topics go deeper into the ideas
%  of the at-level lessons rather than adding more of the same work.
%
%  Topics 1-5 are written here: a reading, the deeper idea worked through
%  (I do / We do), then %%PARTD:N%% pulls in the old Part D questions of
%  at-level Lesson N as practice, then a Challenge. Topics 6 and 7 are the old
%  optional Lessons 11 and 14, pulled in whole by %%LESSON:N%%.
%  Built by make_a5.py. Answer keys: new slots a1.* ... a5.*; reused slots keep
%  their A4 keys (1.x13a ...).
% ══════════════════════════════════════════════════════════════════════════════

% ══════════════════════════════════════════════════════════════════════════════
%  TOPIC 1 — MASS, WEIGHT AND GRAVITY BEYOND EARTH          (deepens Lesson 1)
% ══════════════════════════════════════════════════════════════════════════════
\setcounter{section}{0}\renewcommand{\lessonart}{1}
\sectiondiv{Mass, Weight and Gravity Beyond Earth}{Gravitational field strength \textbullet{} Why $g$ changes from world to world \textbullet{} Balances and spring scales \textbullet{} Weightless in orbit}\label{sec:1}
\lessoncodes{Deepens Lesson 1 \textbullet{} VC2S8U14 \textbullet{} VC2S10U17}

\vocabhead
\begin{ctab}
\begin{tabular}{|SW{4.2cm}|W{11.6cm}|}
  \hd & \thead{Word} & \thead{Definition (copy from the slide)} \\ \hline
  & \scaf{\textbf{Gravitational field strength}} & \scaf{The force of gravity on each kilogram of mass at a place, $g$, in newtons per kilogram (N/kg).} \\ \hline
  & \textbf{Gravitational field} & \cell{a1.v2} \\ \hline
  & \textbf{Free fall} & \cell{a1.v3} \\ \hline
  & \textbf{Orbit} & \cell{a1.v4} \\ \hline
  & \textbf{Apparent weight} & \cell{a1.v5} \\ \hline
\end{tabular}
\end{ctab}

\lessonpart{A}{Reading and comprehension}{Teal}\label{read:1}

\begin{readingbox}{Reading --- Floating in the Space Station}
Astronauts on the International Space Station float about the cabin, so it is natural to think there is no gravity up there. In fact the station orbits only about 400\,km above the ground, and gravity there is still about 90\% as strong as on Earth's surface. If you could build a tower 400\,km tall and stand on top, you would feel almost your normal weight.

So why do astronauts float? Because the station and everything in it are falling. They are moving sideways so fast --- about 7.7\,km every second --- that as they fall towards Earth, the curved ground falls away beneath them just as quickly. They keep falling around the planet without ever getting closer. This endless fall is an \textbf{orbit}, and things falling freely like this are in \textbf{free fall}.

Every planet and moon is surrounded by a \textbf{gravitational field}, the region where its gravity pulls on masses. The \textbf{gravitational field strength}, $g$, tells you how hard it pulls on each kilogram: 9.8\,N/kg on Earth, 1.6\,N/kg on the Moon. A world with more mass pulls harder, but if it is also much wider, its surface is further from its centre, and the pull there is weaker.

What you \emph{feel} as your weight is really the push of the floor holding you up. In free fall nothing pushes on you, so your \textbf{apparent weight} is zero, even though gravity is still pulling.
\end{readingbox}

\rcreset
\rcqk{a1.c1}{How strong is gravity at the height of the Space Station, compared with on Earth's surface?}{1.6cm}
\rcqk{a1.c2}{Explain why astronauts in the Space Station float, even though gravity is still pulling on them.}{2.4cm}
\rcqk{a1.c3}{Name two things about a planet that decide how strong its gravity is at the surface.}{2cm}
\rcqk{a1.c4}{What is the difference between your weight and your apparent weight?}{2.4cm}

\lessonpart{B}{Theory and examples}{Blue}\label{theory:1}

\subsection{Weight in Different Places}\label{sub:a1.1}

\begin{skillbox}{Weight, mass and $g$}
\[ W = m \times g \qquad m = \frac{W}{g} \qquad g = \frac{W}{m} \]
$W$ is weight (N), $m$ is mass (kg), $g$ is gravitational field strength (N/kg). Your mass is the same everywhere. Your weight changes with $g$.
\end{skillbox}

\begin{ctab}
\begin{tabular}{|C{3.0cm}|C{3.8cm}|C{3.8cm}|}
  \hline\thead{Place} & \thead{$g$ (N/kg)} & \thead{Weight of a 50\,kg student (N)} \\ \hline
  Moon & 1.6 & 80 \\ \hline
  Mars & 3.7 & 185 \\ \hline
  Earth & 9.8 & 490 \\ \hline
  Jupiter (cloud tops) & 24.8 & 1240 \\ \hline
\end{tabular}
\end{ctab}

\begin{workedex}{Finding $g$ on an unknown world}
A probe lands on a moon. A 4\,kg test mass hangs from its spring scale, which reads 5.2\,N. Find $g$ on the moon.

\smallskip
$g = W \div m = 5.2 \div 4 = \mathbf{1.3\ \text{N/kg}}$, a little weaker than on our Moon.
\end{workedex}

\begin{wedo}{Finding $g$ on an unknown world}
The same probe's 4\,kg mass reads 15.2\,N on a second world. Find $g$ there, and say which world in the table it is most like.
\wbox{a1.wd1}{3cm}
\end{wedo}

\subsection{Why $g$ Is Different on Different Worlds}\label{sub:a1.2}

\begin{workedex}{Mass and size together}
Jupiter has about 318 times the mass of Earth, but its $g$ is only about 2.5 times Earth's. Why?

\smallskip
More mass means a stronger pull. But Jupiter is also about 11 times as wide as Earth, so its cloud tops are much further from its centre, and gravity weakens with distance. The two effects work against each other, so the extra mass is mostly cancelled out by the extra distance.
\end{workedex}

\begin{wedo}{Mass and size together}
Mars has about one-tenth of Earth's mass and about half of Earth's width. Use both facts to explain why $g$ on Mars (3.7\,N/kg) is less than Earth's, but not as small as one-tenth of it.
\abox{a1.wd2}{4.5cm}
\end{wedo}

\lessonpart{C}{Practice}{Amber}\label{prac:1}

%%PARTD:1%%

\needspace{10\baselineskip}
\task A student stands on bathroom scales in a lift. Explain what happens to the reading as the lift starts to speed up going \emph{up}, and why. (Hint: what must the floor do to make the student speed up?)
\abox{a1.t1}{5cm}

\lessonpart{D}{Challenge}{Violet}\label{ext:1}

\task The Space Station travels at 7.7\,km/s. Explain, step by step, why the astronauts inside it are falling all the time but never land. Then explain what would happen to the station if it suddenly stopped moving sideways.
\abox{a1.ch}{7cm}

\clearpage
% ══════════════════════════════════════════════════════════════════════════════
%  TOPIC 2 — NET FORCE, DRAG AND TERMINAL VELOCITY          (deepens Lesson 2)
% ══════════════════════════════════════════════════════════════════════════════
\setcounter{section}{1}\renewcommand{\lessonart}{2}
\sectiondiv{Net Force, Drag and Terminal Velocity}{Drag grows with speed \textbullet{} Four forces on a vehicle \textbullet{} The skydiver's fall, stage by stage}\label{sec:2}
\lessoncodes{Deepens Lesson 2 \textbullet{} VC2S8U14 \textbullet{} VC2S10U17}

\vocabhead
\begin{ctab}
\begin{tabular}{|SW{4.2cm}|W{11.6cm}|}
  \hd & \thead{Word} & \thead{Definition (copy from the slide)} \\ \hline
  & \scaf{\textbf{Drag}} & \scaf{A force that pushes against an object moving through air or water; air resistance is drag in air.} \\ \hline
  & \textbf{Terminal velocity} & \cell{a2.v2} \\ \hline
  & \textbf{Streamlined} & \cell{a2.v3} \\ \hline
  & \textbf{Driving force} & \cell{a2.v4} \\ \hline
  & \textbf{Resistive force} & \cell{a2.v5} \\ \hline
\end{tabular}
\end{ctab}

\lessonpart{A}{Reading and comprehension}{Teal}\label{read:2}

\begin{readingbox}{Reading --- The Fastest Fall}
In 2012 Felix Baumgartner jumped from a balloon 39\,km above New Mexico. In the thin air up there, almost nothing held him back, and he fell faster than the speed of sound --- over 1300\,km/h. A skydiver jumping from a plane at 4\,km never gets close to that. Why?

The answer is \textbf{drag}. Drag is the force of the air pushing against a moving object, and it grows as the object goes faster. At first a falling skydiver is slow, drag is small, and the net force is almost their whole weight, so they speed up quickly. As they speed up, drag grows, the net force shrinks, and they speed up less and less. Eventually drag is as big as their weight. The forces are balanced, and they fall at a steady top speed: their \textbf{terminal velocity}, about 200\,km/h for a skydiver lying flat.

Thick air means lots of drag. Baumgartner's air was so thin that drag stayed small, so his terminal velocity was huge. Shape matters too: a \textbf{streamlined} shape, like a diving falcon's, slips through the air with less drag.

Cars face the same battle. The engine provides a \textbf{driving force}; drag and friction are \textbf{resistive forces}. A car's top speed is reached when the resistive forces grow to match the driving force.
\end{readingbox}

\rcreset
\rcqk{a2.c1}{What happens to drag as an object moves faster?}{1.6cm}
\rcqk{a2.c2}{Why does a skydiver speed up quickly at first, then less and less?}{2.4cm}
\rcqk{a2.c3}{What is terminal velocity, and when is it reached?}{2.2cm}
\rcqk{a2.c4}{Why could Baumgartner fall so much faster than an ordinary skydiver?}{2.2cm}

\lessonpart{B}{Theory and examples}{Blue}\label{theory:2}

\subsection{The Skydiver, Stage by Stage}\label{sub:a2.1}

\begin{center}
\begin{tikzpicture}[font=\small]
  \draw[-{Stealth[length=5pt]}, line width=0.8pt] (0,0) -- (9.2,0) node[below left, font=\scriptsize]{time};
  \draw[-{Stealth[length=5pt]}, line width=0.8pt] (0,0) -- (0,4.4) node[above right, font=\scriptsize]{speed};
  \draw[line width=1.4pt] (0,0) .. controls (1.2,2.6) and (2.2,3.4) .. (3.6,3.5) -- (5,3.5);
  \draw[line width=1.4pt] (5,3.5) .. controls (5.3,1.6) and (5.6,1.3) .. (6.4,1.2) -- (8.8,1.2);
  \draw[densely dotted] (0,3.5) node[left, font=\scriptsize]{terminal} -- (3.6,3.5);
  \draw[densely dotted] (0,1.2) node[left, font=\scriptsize]{new terminal} -- (6.4,1.2);
  \foreach \x/\y/\l in {0.7/1.55/A, 4.3/3.85/B, 5.6/2.4/C, 7.6/1.55/D}{\node[draw, circle, inner sep=1pt, font=\scriptsize\bfseries, fill=white] at (\x,\y) {\l};}
  \node[font=\scriptsize, anchor=west] at (5.1,3.75) {parachute opens};
  \draw[-{Stealth[length=4pt]}] (5.6,3.6) -- (5.05,3.5);
\end{tikzpicture}
\figcap{Speed of a skydiver: A speeding up, B terminal velocity, C parachute opens, D a new, slower terminal velocity.}
\end{center}

\begin{workedex}{Explaining stage A and stage B}
\textbf{Stage A:} weight (down) is bigger than drag (up). The net force is downwards, so the skydiver speeds up. As speed grows, drag grows, so the net force gets smaller.\\
\textbf{Stage B:} drag has grown until it equals weight. The forces are balanced, the net force is zero, so the speed stays the same: terminal velocity.
\end{workedex}

\begin{wedo}{Explaining stage C and stage D}
Use the words \emph{drag}, \emph{weight} and \emph{net force} to explain what happens to the skydiver's speed at stage C, when the parachute opens, and why it levels off again at stage D.
\abox{a2.wd1}{5cm}
\end{wedo}

\subsection{Four Forces on a Vehicle}\label{sub:a2.2}

\begin{workedex}{A car's top speed}
A car's engine gives a driving force of 3000\,N. At 80\,km/h the resistive forces total 1800\,N.

\smallskip
Net force $= 3000 - 1800 = \mathbf{1200\ \text{N}}$ forwards, so the car is still speeding up. As it goes faster the resistive forces grow. It reaches top speed when they reach 3000\,N, and the net force is zero.
\end{workedex}

\begin{wedo}{A car's top speed}
A cyclist pedals with a driving force of 150\,N. At 30\,km/h the resistive forces are 150\,N. What is the net force, and what is happening to the cyclist's speed? Suggest one way the cyclist could reach a higher top speed without pedalling harder.
\abox{a2.wd2}{4.5cm}
\end{wedo}

\lessonpart{C}{Practice}{Amber}\label{prac:2}

%%PARTD:2%%

\lessonpart{D}{Challenge}{Violet}\label{ext:2}

\task A falcon dives for prey with its wings folded back. A second falcon of the same mass dives with its wings spread out. Predict which reaches the greater terminal velocity, and explain why, using the forces acting on each bird. Sketch both lines on one speed--time graph.
\abox{a2.ch}{5cm}
\drawk{a2.ch-graph}{5cm}{}

\clearpage
% ══════════════════════════════════════════════════════════════════════════════
%  TOPIC 3 — FORCE PAIRS AND MOMENTUM                       (deepens Lesson 4)
% ══════════════════════════════════════════════════════════════════════════════
\setcounter{section}{2}\renewcommand{\lessonart}{4}
\sectiondiv{Force Pairs and Momentum}{Third Law pair or balanced pair? \textbullet{} Momentum \textbullet{} Pushing apart: recoil}\label{sec:3}
\lessoncodes{Deepens Lesson 4 \textbullet{} VC2S8U14 \textbullet{} VC2S10U17}

\vocabhead
\begin{ctab}
\begin{tabular}{|SW{4.2cm}|W{11.6cm}|}
  \hd & \thead{Word} & \thead{Definition (copy from the slide)} \\ \hline
  & \scaf{\textbf{Momentum}} & \scaf{How hard it is to stop a moving object: mass $\times$ velocity, in kg\,m/s.} \\ \hline
  & \textbf{Third Law pair} & \cell{a3.v2} \\ \hline
  & \textbf{Recoil} & \cell{a3.v3} \\ \hline
  & \textbf{Conservation of momentum} & \cell{a3.v4} \\ \hline
\end{tabular}
\end{ctab}

\lessonpart{A}{Reading and comprehension}{Teal}\label{read:3}

\begin{readingbox}{Reading --- The Kick of a Fire Hose}
A fire hose at full pressure can knock a firefighter over, so it takes two or three people to hold one. The hose pushes water forwards out of the nozzle, and by Newton's Third Law the water pushes the hose backwards with an equal force. This backwards push is called \textbf{recoil}. A cannon recoils when it fires, and a skateboard rolls backwards when you jump off its front.

Newton's Third Law pairs always act on two \emph{different} objects --- here, the hose and the water. That is why they do not cancel out. Forces only cancel when they act on the \emph{same} object.

Scientists use \textbf{momentum} to describe recoil. Momentum is mass times velocity: a heavy, fast object has a lot of momentum. When two things push apart from rest, they end up with equal momentum in opposite directions, so the total is still zero. This is the law of \textbf{conservation of momentum}. A small, light object shoots away fast; a big, heavy one moves off slowly.
\end{readingbox}

\rcreset
\rcqk{a3.c1}{Why does a fire hose push backwards on the firefighters?}{2.2cm}
\rcqk{a3.c2}{Why do the two forces in a Third Law pair not cancel each other out?}{2.2cm}
\rcqk{a3.c3}{What two things decide how much momentum an object has?}{1.6cm}
\rcqk{a3.c4}{Two objects push apart from rest. Which one moves faster, the light one or the heavy one?}{1.8cm}

\lessonpart{B}{Theory and examples}{Blue}\label{theory:3}

\subsection{Third Law Pair or Balanced Pair?}\label{sub:a3.1}

\begin{skillbox}{The three tests for a Third Law pair}
\textbf{1.} The two forces act on \emph{two different objects}. \quad \textbf{2.} They are the \emph{same type} of force (both gravity, both contact \dots). \quad \textbf{3.} They are equal in size and opposite in direction.

A \textbf{balanced pair} passes only test 3: both forces act on the same object.
\end{skillbox}

\begin{workedex}{A book on a table}
\textbf{Earth pulls the book down (weight)} and \textbf{the table pushes the book up (normal force)}: both act on the book, and they are different types. This is a \emph{balanced pair}, not a Third Law pair.

\smallskip
The Third Law partner of the book's weight is \textbf{the book pulling the Earth up} --- same type (gravity), on a different object (the Earth).
\end{workedex}

\begin{wedo}{A swimmer}
A swimmer pushes backwards on the water with her hands. Name the Third Law partner of this force, and use the three tests to show it is a Third Law pair.
\abox{a3.wd1}{4.5cm}
\end{wedo}

\subsection{Momentum and Recoil}\label{sub:a3.2}

\begin{skillbox}{Momentum}
\[ p = m \times v \]
$p$ is momentum (kg\,m/s), $m$ is mass (kg), $v$ is velocity (m/s). When two objects push apart from rest: \quad $m_1 v_1 = m_2 v_2$ (in opposite directions).
\end{skillbox}

\begin{workedex}{Two skaters}
A 60\,kg skater and a 40\,kg skater stand still on ice and push apart. The 60\,kg skater moves off at 2\,m/s. How fast does the 40\,kg skater move?

\smallskip
Momentum of the first skater $= 60 \times 2 = 120$\,kg\,m/s.\\
The second must have 120\,kg\,m/s the other way: $v = 120 \div 40 = \mathbf{3\ \text{m/s}}$.
\end{workedex}

\begin{wedo}{A cannon}
A 500\,kg cannon fires a 10\,kg ball at 100\,m/s. Find the ball's momentum, then the speed at which the cannon recoils.
\wbox{a3.wd2}{3.5cm}
\end{wedo}

\lessonpart{C}{Practice}{Amber}\label{prac:3}

%%PARTD:4%%

\lessonpart{D}{Challenge}{Violet}\label{ext:3}

\task A rocket in deep space has nothing to push against, yet it can still speed up. Use Newton's Third Law \emph{and} the idea of momentum to explain how. Then explain why a rocket that throws its exhaust out faster gains more speed from the same amount of fuel.
\abox{a3.ch}{7cm}

\clearpage
% ══════════════════════════════════════════════════════════════════════════════
%  TOPIC 4 — F = ma: ACCELERATION AND IMPACT                (deepens Lesson 5)
% ══════════════════════════════════════════════════════════════════════════════
\setcounter{section}{3}\renewcommand{\lessonart}{5}
\sectiondiv{\texorpdfstring{$F = ma$}{F = ma}: Acceleration and Impact}{Net force first \textbullet{} Acceleration from speeds and time \textbullet{} Lifts \textbullet{} Impacts and crumple zones}\label{sec:4}
\lessoncodes{Deepens Lesson 5 \textbullet{} VC2S8U14 \textbullet{} VC2S10U17}

\vocabhead
\begin{ctab}
\begin{tabular}{|SW{4.2cm}|W{11.6cm}|}
  \hd & \thead{Word} & \thead{Definition (copy from the slide)} \\ \hline
  & \scaf{\textbf{Change in velocity}} & \scaf{Final velocity minus starting velocity, $\Delta v$, in m/s.} \\ \hline
  & \textbf{Deceleration} & \cell{a4.v2} \\ \hline
  & \textbf{Impact force} & \cell{a4.v3} \\ \hline
  & \textbf{Crumple zone} & \cell{a4.v4} \\ \hline
\end{tabular}
\end{ctab}

\lessonpart{A}{Reading and comprehension}{Teal}\label{read:4}

\begin{readingbox}{Reading --- Cars Built to Crush}
A modern car is designed to be destroyed in a crash. Its front and back are \textbf{crumple zones}: parts built to fold up slowly, like a drink can being squashed. This might sound like bad engineering, but it saves lives.

In a crash, the car and the people in it must lose all their speed. Newton's Second Law says that the force needed to change an object's motion depends on how quickly the change happens. If a car stops dead against a wall in a few thousandths of a second, the people inside feel an enormous \textbf{impact force}. If the front crumples and the stop takes ten times as long, the force on them is ten times smaller.

Slowing down is called \textbf{deceleration}. Seatbelts stretch a little, airbags squash, and crumple zones fold, all for the same reason: to make the deceleration take longer, so the force is smaller. Bicycle helmets, the sand pit in long jump and the bend in your knees when you land from a jump all work the same way.
\end{readingbox}

\rcreset
\rcqk{a4.c1}{What is a crumple zone, and why is it built to crush?}{2.2cm}
\rcqk{a4.c2}{How does making a stop take longer change the force on the people in the car?}{2.2cm}
\rcqk{a4.c3}{Give two other examples that work in the same way as a crumple zone.}{2cm}

\lessonpart{B}{Theory and examples}{Blue}\label{theory:4}

\subsection{Acceleration From Speeds and Time}\label{sub:a4.1}

\begin{skillbox}{Two steps}
\textbf{Step 1:} \quad $a = \dfrac{\Delta v}{t} = \dfrac{\text{final velocity} - \text{starting velocity}}{\text{time}}$ \qquad \textbf{Step 2:} \quad $F = m \times a$ \quad (this $F$ is the \emph{net} force).

A negative acceleration is a deceleration.
\end{skillbox}

\begin{workedex}{A car braking}
A 1000\,kg car slows from 20\,m/s to 0\,m/s in 4\,s.

\smallskip
$a = (0 - 20) \div 4 = -5$\,m/s$^2$ \quad (a deceleration of 5\,m/s$^2$)\\
$F = 1000 \times 5 = \mathbf{5000\ \text{N}}$ backwards, from the brakes and road.
\end{workedex}

\begin{wedo}{A sprinter}
A 60\,kg sprinter speeds up from 0 to 9\,m/s in 1.5\,s. Find her acceleration, then the net force on her.
\wbox{a4.wd1}{3.5cm}
\end{wedo}

\subsection{Impacts: the Same Change, a Longer Time}\label{sub:a4.2}

\begin{workedex}{Crash test}
A 70\,kg crash-test dummy stops from 15\,m/s. Compare a stop in 0.01\,s (rigid car) with a stop in 0.1\,s (crumple zone and airbag).

\smallskip
\textbf{Rigid:} $a = 15 \div 0.01 = 1500$\,m/s$^2$; $F = 70 \times 1500 = \mathbf{105\,000\ \text{N}}$\\
\textbf{Crumple:} $a = 15 \div 0.1 = 150$\,m/s$^2$; $F = 70 \times 150 = \mathbf{10\,500\ \text{N}}$, ten times smaller.
\end{workedex}

\begin{wedo}{Landing from a jump}
A 50\,kg student jumps off a wall and hits the ground at 4\,m/s. Find the force on her legs if she lands stiff-legged and stops in 0.02\,s, and if she bends her knees and stops in 0.2\,s.
\wbox{a4.wd2}{4.5cm}
\end{wedo}

\lessonpart{C}{Practice}{Amber}\label{prac:4}

%%PARTD:5%%

\lessonpart{D}{Challenge}{Violet}\label{ext:4}

\task You are designing packaging to post a 0.4\,kg glass vase. In a drop, it hits the ground at 5\,m/s, and it breaks if the force on it is more than 200\,N.
\pt{a}What is the shortest stopping time that keeps the force below 200\,N?
\wbox{a4.ch-a}{3cm}
\pt{b}Describe packaging that would make the stop take at least that long, and explain how it works.
\abox{a4.ch-b}{4.5cm}

\clearpage
% ══════════════════════════════════════════════════════════════════════════════
%  TOPIC 5 — CALCULATING ENERGY                             (deepens Lesson 7)
% ══════════════════════════════════════════════════════════════════════════════
\setcounter{section}{4}\renewcommand{\lessonart}{7}
\sectiondiv{Calculating Energy}{Gravitational potential energy \textbullet{} Kinetic energy \textbullet{} Energy is conserved \textbullet{} Speed from height}\label{sec:5}
\lessoncodes{Deepens Lesson 7 \textbullet{} VC2S8U15 \textbullet{} VC2S10U15}

\vocabhead
\begin{ctab}
\begin{tabular}{|SW{4.2cm}|W{11.6cm}|}
  \hd & \thead{Word} & \thead{Definition (copy from the slide)} \\ \hline
  & \scaf{\textbf{Joule}} & \scaf{The unit of energy, J. Lifting a 1\,N weight up 1\,m takes 1\,J.} \\ \hline
  & \textbf{Law of Conservation of Energy} & \cell{a5.v2} \\ \hline
  & \textbf{Squared} & \cell{a5.v3} \\ \hline
  & \textbf{Sankey diagram} & \cell{a5.v4} \\ \hline
\end{tabular}
\end{ctab}

\lessonpart{A}{Reading and comprehension}{Teal}\label{read:5}

\begin{readingbox}{Reading --- Why Speed Limits Matter}
A car at 60\,km/h and a car at 50\,km/h do not seem very different. But if both brake hard, the faster car needs about 44\% more distance to stop, and if it hits something, it does so with almost half as much energy again. A small rise in speed makes a big rise in energy.

The reason is in the formula for kinetic energy: $\text{KE} = \tfrac{1}{2} m v^2$. The speed is \textbf{squared}, multiplied by itself. Double the speed and the kinetic energy becomes \emph{four} times as big; triple it and it becomes nine times as big. Brakes have to take all that energy away, turning it into thermal energy in the brake pads, so stopping distance grows the same way.

Energy is never made or destroyed, only changed from one form to another: the \textbf{Law of Conservation of Energy}. When a ball falls, its gravitational potential energy, $\text{GPE} = mgh$, turns into kinetic energy. If no energy is wasted, every joule of GPE lost becomes a joule of KE gained --- which lets you work out how fast something will be going at the bottom of a fall or a ramp.
\end{readingbox}

\rcreset
\rcqk{a5.c1}{What happens to a car's kinetic energy when its speed doubles?}{1.6cm}
\rcqk{a5.c2}{Why does a faster car need a much longer distance to stop?}{2.2cm}
\rcqk{a5.c3}{State the Law of Conservation of Energy in your own words.}{2cm}

\lessonpart{B}{Theory and examples}{Blue}\label{theory:5}

\subsection{Two Formulas}\label{sub:a5.1}

\begin{skillbox}{GPE and KE}
\[ \text{GPE} = m \times g \times h \qquad\qquad \text{KE} = \tfrac{1}{2} \times m \times v^2 \]
$m$ in kg, $g = 9.8$\,N/kg, $h$ in m, $v$ in m/s; energy in joules (J). Square the speed \emph{first}, then multiply.
\end{skillbox}

\begin{workedex}{Kinetic energy}
A 0.06\,kg tennis ball is served at 50\,m/s.

\smallskip
$v^2 = 50 \times 50 = 2500$ \qquad $\text{KE} = \tfrac{1}{2} \times 0.06 \times 2500 = \mathbf{75\ \text{J}}$
\end{workedex}

\begin{wedo}{Kinetic energy}
Find the kinetic energy of a 70\,kg cyclist and bike moving at 10\,m/s, and again at 20\,m/s. How many times bigger is the second answer?
\wbox{a5.wd1}{4cm}
\end{wedo}

\subsection{Speed From Height}\label{sub:a5.2}

\begin{workedex}{A skateboarder on a ramp}
A 50\,kg skateboarder rolls from rest down a ramp 1.8\,m high. Ignoring friction, how fast is she going at the bottom?

\smallskip
GPE at the top $= 50 \times 9.8 \times 1.8 = 882$\,J. All of it becomes KE at the bottom.\\
$882 = \tfrac{1}{2} \times 50 \times v^2$, so $v^2 = 882 \div 25 = 35.3$, and $v = \sqrt{35.3} \approx \mathbf{5.9\ \text{m/s}}$.
\end{workedex}

\begin{wedo}{A dropped ball}
A 0.2\,kg ball is dropped from 5\,m. Ignoring air resistance, find its speed just before it lands. Does the mass of the ball matter? Explain.
\wbox{a5.wd2}{5cm}
\end{wedo}

\lessonpart{C}{Practice}{Amber}\label{prac:5}

%%PARTD:7%%

\lessonpart{D}{Challenge}{Violet}\label{ext:5}

\task A roller coaster's first hill is 40\,m high, and a loop further along is 25\,m high. The train starts at rest at the top of the first hill.
\pt{a}Ignoring friction, find the train's speed at the bottom of the first hill and at the top of the loop. (Use energy per kilogram: you do not need the train's mass.)
\wbox{a5.ch-a}{5cm}
\pt{b}In reality 15\% of the energy is wasted before the loop. Will the train still get over a 35\,m hill placed after the loop? Show your reasoning.
\wbox{a5.ch-b}{5cm}

\clearpage
% ══════════════════════════════════════════════════════════════════════════════
%  TOPIC 6 — TURNING EFFECTS AND THE TREBUCHET   (old optional Lesson 11)
% ══════════════════════════════════════════════════════════════════════════════
\setcounter{section}{5}\renewcommand{\lessonart}{11}
\sectiondiv{Turning Effects and the Trebuchet}{Turning effects \textbullet{} Arcs and tip speed \textbullet{} Energy and launch speed \textbullet{} Deepens Lessons 9 and 10}\label{sec:6}
\lessoncodes{Deepens Lessons 9 and 10 \textbullet{} VC2S8U13 \textbullet{} VC2S8U15 \textbullet{} VC2S10U15}
%%LESSON:11%%

% ══════════════════════════════════════════════════════════════════════════════
%  TOPIC 7 — OPTIMISING RANGE                    (old optional Lesson 14)
% ══════════════════════════════════════════════════════════════════════════════
\setcounter{section}{6}\renewcommand{\lessonart}{14}
\sectiondiv{Optimising Range}{Launching at an angle \textbullet{} The best angle \textbullet{} Real launches \textbullet{} Deepens Lesson 12}\label{sec:7}
\lessoncodes{Deepens Lesson 12 \textbullet{} VC2S8U14 \textbullet{} VC2S10U17}
%%LESSON:14%%
